\subsection{同类型分次代数的张量积}

在第 7 号命题 10 的假设成立的前提下，进一步设所有 \(\Delta_{i}\) 都等于同一个交换幺半群 \(\Delta_{0}\) ；于是我们可以在张量 \(\varepsilon\) -积 \({}^{\varepsilon}\bigotimes_{i\in I}E_{i}\) 上考虑类型为 \(\Delta_{0}\) 的总分次，它与此代数上类型为 \(\Delta = \Delta_{0}^{I}\) 的分次相关联（II, § 11, 第 1 号）；我们将带有这个分次的 \({}^{\varepsilon}\bigotimes_{i\in I}E_{i}\) 称为一个分次张量 \(\varepsilon\) -积，其类型为 \(\Delta_{0}\) ，所属的族为 \((E_{i})_{i\in I}\) （族中各代数的类型为 \(\Delta_{0}\) ）.

保持第 7 号命题 10 的记号，设 F 也是一个类型为 \(\Delta_0\) 的分次 A-代数，且诸 \(f_i\) 是类型为 \(\Delta_0\) 的分次代数同态。则 \(f: {}^{\varepsilon}\bigotimes_{i \in I} E_i \to F\) 也是类型为 \(\Delta_0\) 的分次代数同态：因为由公式 (30)（第 7 号）可知，若 \(x_i\) 是齐次的且次数为 \(\alpha_i \in \Delta_0\) ，则 \(\bigotimes_{i \in I} x_i\) 与 \(\prod_{i \in I} f_i(x_i)\) 都是次数为 \(\sum_{i \in I} \alpha_i \in \Delta_0\) 的齐次元素.

因此可以说， \(\bigotimes_{i\in I}E_i\) （带有类型为 \(\Delta_0\) 的总分次）与典范映射 \(\psi\) 一起构成第三个泛映射问题的一个解，其中 \(\Sigma\) 是类型为 \(\Delta_0\) 的分次 A-代数的种类，态射是类型为 \(\Delta_0\) 的分次代数同态，而 \(\alpha\) -映射是映射 \(\prod_{i}f_{i}\) （其中除条件 (26)（第 7 号）外，还假设每个 \(f_{i}\) 是类型为 \(\Delta_0\) 的分次代数同态）.

对 I 的每个子集 J，典范同态 \({}^{\varepsilon}\bigotimes_{i\in J}E_i\to {}^{\varepsilon}\bigotimes_{i\in I}E_i\) （第 7 号）实际上是一个类型为 \(\Delta_0\) 的分次代数同态，这由上述内容立即得出。

命题 12.（同类型分次代数的张量 \(\varepsilon\) -积的“结合性”）。采用第 7 号命题 10 的记号，设所有 \(\Delta_{i}\) 都等于同一个幺半群 \(\Delta_{0}\) ；设 \((J_{\lambda})_{\lambda \in L}\) 是 I 的一个划分。采用第 7 号命题 11 的记号，设公式 (32)（第 7 号）的右端只依赖于和 \(\alpha_{\lambda}^{\prime \prime} = \sum_{i \in J_{\lambda}} \alpha_{i}, \beta_{\mu}^{\prime \prime} = \sum_{j \in J_{\mu}} \beta_{j}\) ——对相异指标的每个有序对 \((\lambda, \mu)\) 、所有 \(\alpha_{\lambda}^{\prime} \in \Delta_{\lambda}^{\prime}\) 以及所有 \(\beta_{\mu}^{\prime} \in \Delta_{\mu}^{\prime}\) 而言；令 \(\varepsilon_{\lambda \mu}^{\prime \prime}(\alpha_{\lambda}^{\prime \prime}, \beta_{\mu}^{\prime \prime})\) 表示 (32) 的右端。则 \((\varepsilon_{\lambda \mu}^{\prime \prime})\) 是族 \((\Delta_{\lambda}^{\prime \prime})_{\lambda \in L}\) 上的一个交换因子系，其中 \(\Delta_{\lambda}^{\prime \prime} = \Delta_{0}\) （对所有 \(\lambda \in L\) ）。若 \(E_{\lambda}^{\prime \prime}\) 是一个分次张量 \(\varepsilon\) -积，其类型为 \(\Delta_{0}\) ，所属的族为 \((E_{i})_{i \in J_{\lambda}}\) ，则存在且仅存在一个类型为 \(\Delta_{0}\) 的分次代数同构 \(w: {}^{\varepsilon}\bigotimes_{\lambda \in L} E_{\lambda}^{\prime \prime} \to {}^{\varepsilon}\bigotimes_{i \in I} E_{i}\) ，使得

\[
w \left(\bigotimes_ {\lambda \in L} \left(\bigotimes_ {i \in J _ {\lambda}} x _ {i}\right)\right) = \bigotimes_ {i \in I} x _ {i}\tag{35}
\]

前提是按第 7 号命题 11 所述，在 \(\mathbf{J}_{\lambda}\) 与 \(\mathbf{I}\) 上选取全序。

根据假设，对 \(\gamma\) , \(\delta\) （属于 \(\Delta_0\) ），有 \(\varepsilon_{\lambda \mu}''(\gamma, \delta) = \varepsilon_{i_0 j_0}(\gamma, \delta)\) （对某个 \(i_0 \in J_\lambda\) 与某个 \(j_0 \in J_\mu\) ），这可以通过考虑元素 \(\alpha_\lambda' = (\alpha_i)_{i \in J_\lambda}\) 与 \(\beta_\mu' = (\beta_j)_{j \in J_\mu}\) 看出，它们满足 \(\alpha_{i_0} = \gamma\) , \(\alpha_i = 0\) （对 \(i \neq i_0\) ）, \(\beta_{j_0} = \delta\) , \(\beta_j = 0\) （对 \(j \neq j_0\) ）；由此立即得出，诸 \(\varepsilon_{\lambda \mu}''\) 构成一个交换因子系。证明的其余部分与命题 11（第 7 号）的证明类似，留给读者完成。

注意，当 \(\Delta_0 = \mathbf{Z}\) 且 \((\varepsilon_{ij})\) 是平凡的因子系 \((\varepsilon_{ij}(\alpha_i,\beta_j) = 1\) 对所有 \(i,j)\) ，或是由 \(\varepsilon_{ij}(\alpha_i,\beta_j) = (-1)^{\alpha_i\beta_j}\) 定义的因子系时，命题 12 的附加假设成立；在后一情形，公式 (32) 的右端等于 \((-1)^{\gamma}\) ，其中 \(\gamma = \sum_{i\in J_{\lambda},j\in J_{\mu}}\alpha_{i}\beta_{j} = \left(\sum_{i\in J_{\lambda}}\alpha_{i}\right)\left(\sum_{j\in J_{\mu}}\beta_{j}\right)\) .

注。(1) 设 I 是一个无限指标集， \(\Delta_0\) 是一个交换幺半群；设 \((\Delta_i)_{i \in I}\) 表示满足 \(\Delta_i = \Delta_0\) （对所有 \(i\) ）的族，并对 I 的每对相异指标 \((i,j)\) 给定一个满足条件 (22)、(23) 与 (24)（第 7 号）的映射 \(\varepsilon_{ij} : \Delta_i \times \Delta_j \to A\) ；这也称为族 \((\Delta_i)\) 上的一个交换因子系。考虑一个族 \((E_i)_{i \in I}\) （由类型为 \(\Delta_0\) 的分次 A-代数组成）；对 I 的每个有限子集 J，令 \(E_J\) 表示一个分次张量 \(\varepsilon\) -积，其类型为 \(\Delta_0\) ，所属的子族为 \((E_i)_{i \in J}\) （对 J 上的全序作任意选取）。若 J 与 \(J'\) 是 I 的两个有限子集且 \(J \subset J'\) ，则上面已定义了一个类型为 \(\Delta_0\) 的分次代数的典范同态 \(h_{J^{\prime}J}:E_{J}\to E_{J^{\prime}}\) ，而这些同态的唯一性性质立即表明：若 \(J\subset J^{\prime}\subset J^{\prime\prime}\) 是 I 的三个有限子集，则 \(h_{J^{\prime\prime}J}=h_{J^{\prime\prime}J^{\prime}}\circ h_{J^{\prime}J}\) 。于是有一个正向系统 \((E_{J},h_{J^{\prime}J})\) ——由类型为 \(\Delta_{0}\) 的分次代数组成（§3, 第 3 号）——其指标集是 I 的有限子集组成的右有向集 \(\mathfrak{F}(I)\) 。类型为 \(\Delta_{0}\) 的分次代数——这个正向系统的正向极限（§3, 第 3 号）——称为一个分次张量 \(\varepsilon\) -积，其类型为 \(\Delta_{0}\) ，所属的族为 \((E_{i})_{i\in I}\) ；它也记作 \({}^{\varepsilon}\bigotimes_{i\in I}E_{i}\) 。当所有 \(\Delta_{i}\) 都等于 \(\mathbf{Z}\) 且 \(\varepsilon_{ij}(\alpha_{i},\beta_{j})=(-1)^{\alpha_{i}\beta_{j}}\) 时，张量积 \({}^{\varepsilon}\bigotimes_{i\in I}E_{i}\) 也称为

族 \((\mathbf{E}_{i})_{i\in\mathbb{I}}\) 的斜张量积，记作 \({}^{g}\bigotimes_{i\in I}E_{i}\) 。我们把如下命题的表述与证明留给读者：该命题把第 7 号的命题 10 推广到 I 为无限的情形，正如第 5 号的命题 8 把第 2 号的命题 5 推广到 I 为无限的情形那样。注意， \({}^{\varepsilon}\bigotimes_{i\in I}E_{i}\) 的底 A-模与第 5 号中定义的非分次代数的族 \((\mathbf{E}_{i})_{i\in\mathbb{I}}\) 的（非分次）张量积的底 A-模相同。

\begin{enumerate}
\def\labelenumi{(\arabic{enumi})}
\setcounter{enumi}{1}
\tightlist
\item
  设 E 是一个类型为 \(\Delta_0\) 的分次 A-代数（其中 \(\Delta_0\) 是一个交换幺半群），且 \(\rho: \mathbf{A} \to \mathbf{B}\) 是一个环同态； \(\rho^*(\mathbf{E})\) 上的分次（II, §11, 第 5 号）与分次张量积 \(\mathbf{B} \otimes_{\mathbf{A}} \mathbf{E}\) （其中 B 带有平凡分次）上的分次相同。
\end{enumerate}
% semantic-fragments: legacy-input-seam
\relax{}
